\documentclass[openany,fontset=fandol]{ctexbook}
\usepackage{amsmath,amssymb}
\usepackage[unicode,hidelinks]{hyperref}
\setlength{\emergencystretch}{2em}
\interfootnotelinepenalty=10000
\DeclareMathOperator{\Spec}{Spec}
\DeclareMathOperator{\Aut}{Aut}
\DeclareMathOperator{\Div}{Div}
\newcommand{\OO}{\mathcal O}
\newcommand{\PP}{\mathbf P}
\newcommand{\m}{\mathfrak m}
\newcommand{\Hom}{\operatorname{Hom}}
\raggedbottom
\begin{document}
\title{代数几何基础第五卷预备笔记\\任意底上的单态性与无穷小差}
\author{据法文扫描和英译编辑稿互校续译}
\date{2026年11月22日}
\maketitle
\chapter*{范围与接续}
第84批续译第三篇命题4.3的任意底证明及其后关于参数丛、下降的作者说明。
法文物理143--146页（手写150--153），英译物理108--110页。
译文从“现在概述一般情形”起，到扭曲射影丛的说明止。
下一处是法文物理146页、英译物理110页“再回到较为具体的情形”，随即进入命题4.4。
本批不改动正式出版的 EGA IV.1，也未并入1158页合订本。

\chapter*{第三篇（续）}
\section*{4\quad 除子线性系与可逆模（续）}
\noindent\textbf{命题4.3的证明：任意底。}
现在概述一般情形，即 \(S\) 任意。换底之后，a) 只须证明：\(P^\vee\) 在 \(S\) 上定义了 \(X\) 上同一个除子的两条截面必相等；b) 只须证明：满足 \(D\circ u=D\circ v\) 的 \(P^\vee\) 的两个自同构必相等。
可以假定 \(S\) 仿射。在 b) 中，没有专门假定 \(g\) 有限表示，但 \(g\) 满射。利用 \(g\) 是开映射，可以局部归结到 \(X\) 也仿射，因而在 \(S\) 上有限表示的情形。再用惯用的逼近方法，归结到 \(S\) 诺特的情形。\footnote{法143的“可以假定 \(S\) 仿射”为手写补入。b) 归约时，取 \(X\) 的仿射开，其像覆盖底；围绕底点选入该像内的主开，并限制该仿射开到这个主开。不能任取一个仿射开就断言它仍满射到底。原稿把诺特逼近留作惯用步骤；下文补证明确比较的有限数据。}

设底概形 \(S\) 诺特，且有 \(S\) 上的函子态射 \(h:F\to G\)；\(F,G\) 都是从 \((\mathrm{Sch}/S)^{\mathrm{op}}\) 到集合范畴的函子。有很一般的判据（将在第五章阐明）：在关于 \(F,G\) 的一些简单假设下，如果每个 \(s\in S\) 上的纤维态射 \(F_s\to G_s\) 都是单态射，则 \(F\to G\) 是单态射。这里记 \(F_s=F\times_S\Spec k(s)\)，\(G_s\) 同理。举例说，当两者都由 \(S\) 上局部有限型的概形表出时，这些假设成立；但在我们这里，先验有表出的只有第一个函子。现就所关心的两个特例概述第五章的论证\nobreak\footnote{原稿说的是纤维函子的单态射，也就是对纤维上的所有检验概形单射，而不仅是剩余域点的单射。英108漏写了法143例子中的“局部”。这里的第五章是手稿中预告的章节。}。

我们有有限型 \(S\)-概形 \(F\) 的两个截面 \(u,v\)（分别取 \(F=P^\vee\) 或射影群 \(F=\mathrm{GP}(E^\vee)\)），要证明二者相等。只须对每个 \(s\in S\) 和每个 \(n\geq0\)，在换底
\[
 \Spec\bigl(\OO_{S,s}/\m_s^{n+1}\bigr)\longrightarrow S
\]
后证明两截面相等。\footnote{法143与英108一笔带过，此处由闭对角线形成等化子，再用克鲁尔交定理；见补证。}
于是归结到阿廷局部底 \(S=\Spec A\)。设 \(\m^{n+1}=0\)，按 \(n\) 归纳。可假定 \(u,v\) 模 \(\m^n\) 相等，于是它们相差一个元素
\[
 \delta\in\Hom_k\bigl(u_0^*\Omega^1_{F_0/k},V\bigr).
\]
这里 \(k=A/\m\) 是剩余域，\(F_0=F\times_A\Spec k\) 是剩余域纤维，\(V=\m^n=\m^n/\m^{n+1}\)。利用 \(h(u)=h(v)\)，只须证明 \(\delta=0\)。\footnote{法144、英109称 \(F_0\) 为“约化纤维”。公式只对底取剩余域，并不把所得纤维再取约化，否则可能丢失切向量。归纳的 \(n=0\) 步直接用纤维单态射；对 \(n\ge1\)，\(V^2=0\) 且 \(\m V=0\)。}

验证的一般原则如下。先把 \(h(u),h(v)\) 看成模 \(\m^n\) 一致的两个点。它们的差可写成
\[
 \delta'\in\Hom_k\bigl(w_0^*\Omega^1_{G_0/k},V\bigr),
 \qquad w_0=h_0(u_0)=h_0(v_0),\quad G_0=G\times_A\Spec k.
\]
这个元素恰由 \(\delta\) 与 \(h_0:F_0\to G_0\) 所诱导的自然映射
\[
 h_0^*:w_0^*\Omega^1_{G_0/k}\longrightarrow u_0^*\Omega^1_{F_0/k}
\]
复合得到。因为 \(h(u)=h(v)\)，\(\delta'=0\)，即 \(\delta\) 与 \(h_0^*\) 的复合为零。因此只要 \(h_0^*\) 满射，便有 \(\delta=0\)。\(h_0\) 是单态射，特别是在 \(k\) 的对偶数环值点上诱导单射，所以 \(h_0^*\) 的对偶单射，\(h_0^*\) 也就满射。\footnote{法144的中间一句漏写星号，字面写成“\(h_0\) 满射”；所需的是前一行的余切映射 \(h_0^*\) 满射，而不是函子态射满射。英109同样有此漏写。}

这一论证可原样用于可表出的 \(G\)，但所关心的情形不属于此类。尽管如此，可定义一个 \(k\)-向量空间 \(T_{w_0}G_0\)，充当 \(w_0^*\Omega^1_{G_0/k}\) 的对偶所扮演的角色，即 \(G_0\) 在 \(w_0\) 的切空间。把两个模 \(\m^n\) 一致的 \(G\) 的点的偏差，表示为 \(T_{w_0}G_0\otimes_kV\) 的元素。这是直接的计算，已包含在手稿拟援引的无穷小扩张部分。\footnote{法144手写补入“切空间”及张量积。此处节号有涂改，英109暂作§26并带问号；下一页作者说明所引“小计算”的节号按英110读作§25。两处均属手稿内部预告，不改作正式出版章节。}

在 a) 中，\(G\) 是 \(\Div_{X/S}\)。于是 \(w_0\) 对应于 \(X_0=X\times_A\Spec k\) 上的 Cartier 除子 \(D_0\)，所取的切空间是
\[
 T_{w_0}G_0=H^0(D_0,\mathcal N_{D_0/X_0}),
 \qquad
 \mathcal N_{D_0/X_0}\simeq\OO_{X_0}(D_0)|_{D_0}.
\]
这里 \(\mathcal N_{D_0/X_0}\) 是 \(D_0\) 在 \(X_0\) 中的法丛。

在 b) 中，可以把 \(X\) 限制到共同基点集的补开，因而设 \(D\) 无基点。更便于用到 \(P\) 的态射来理解此情形（见4.1）。这时 \(G\) 是相对态射函子 \(\underline{\Hom}_S(X,P)\)，而 \(T_{w_0}G_0\) 是\footnote{法145、英109把 Hom 的靶写成 \(\OO_{P_0}\)。左项已拉回到 \(X_0\)，差分是到 \(X_0\) 上结构层的导子，故正文改为 \(\OO_{X_0}\)。取基点补开并不要求一般底上的仿射开本身全局满射。}
\[
 \Hom_{\OO_{X_0}}\bigl(f_0^*\Omega^1_{P_0/k},\OO_{X_0}\bigr).
\]


在两个情形中，都有自然线性映射
\[
 T_{u_0}F_0\otimes_kV\longrightarrow T_{w_0}G_0\otimes_kV,
 \qquad T_{u_0}F_0=\bigl(u_0^*\Omega^1_{F_0/k}\bigr)^\vee .
\]
它表示应用 \(h\) 后，偏差 \(\delta\) 变为对应的 \(\delta'\)。
该映射是单射，因为 \(T_{u_0}F_0\to T_{w_0}G_0\) 是单射；后者来自 \(h_0:F_0\to G_0\) 是单态射。

\medskip\noindent\textbf{作者的写作说明。}
实际撰写时，证明的最后一部分似乎难以不引用第25节的几个小计算；没有必要在此特例中重作这些计算。应指出，这不会构成循环论证，因为第25节及其中发展的计算只依赖第16节的微分演算结果；此外，在第四章中将不再使用4.3，至多可能在接下来的两小节使用。\footnote{这一段在法145左侧有斜线，似为编辑标记。英110编辑者猜测这部分原拟作 EGA 第四章的片段；这只是编排猜测，不属于命题。}

4.3 a) 的意义在于，在所述条件下，参数射影丛可解释为 \(\Div_{X/S}\) 的一个子函子。这样便可以内在地解释 \(X/S\) 上的线性系的同构类（同构允许改变参数射影丛）：它是 \(\Div_{X/S}\) 的一个子函子 \(P^\vee\)，由射影丛表出，而且由其典范包含所定义的除子族在第3小节的意义下为线性。这实质上是经典观点：除子线性系就是满足一定条件的一组除子（比较4.4）。\footnote{法145恢复两处 \(\Div_{X/S}\)。文中的“改变参数丛”指在同构携带除子族的意义下，并不忘记两个参数丛间的相容性。法稿末尾节号有覆写，英110取3；这里的“线性”沿用前文定义。}

另一方面，在 \(g\) 满射时，4.3 b) 也就是说，如果分别以两个射影丛 \(P^\vee,(P^\vee)'\) 参数化的 \(X/S\) 上的两个除子线性系同构，则存在唯一的 \(P^\vee\to(P^\vee)'\) 的同构，与 \(D,D'\) 相容。因此，可以说 \(X/S\) 上线性系的一个同构类确定其参数射影丛，精确到唯一同构。在技术上，等第五章的下降理论可用时，就能据此对除子线性系作忠实平坦下降；不过参数丛还须容许“扭曲射影丛”，这将在后续小节处理。\footnote{法146手写补入“当 \(g\) 满射时”，英110保留。原文在这里预告下降，并未证明全面的可表出性或下降定理。}

\section*{译校补证（非原文）}
\paragraph{从所有阶厚化回到诺特底。}
本命题中的 \(F\) 在底上分离。把其对角线沿 \((u,v):S\to F\times_SF\) 拉回，得到由理想层 \(\mathcal I\) 定义的截面等化子。
若所有阶厚化上相等，则 \(\mathcal I_s\subseteq\bigcap_{n\geq0}\m_s^{n+1}\OO_{S,s}=0\)（\href{https://stacks.math.columbia.edu/tag/00IP}{Stacks，00IP}）。这对每个 \(s\) 都成立，故 \(\mathcal I=0\)。

\paragraph{两个提升的差。}
设 \(n\geq1\)。在两截面具有共同像的仿射邻域，两个环映射在模 \(\m^n\) 后相同，差落在 \(V\)。差对乘积满足导子规则，因为两个差的积落在 \(V^2=0\)，而 \(V\) 上的乘法作用经剩余域。它等价于一个从 \(u_0^*\Omega^1_{F_0/k}\) 到 \(V\) 的线性映射。差为零恰表示两提升相等。复合态射时对导子应用链式法则。

对于相对有效 Cartier 除子，局部取两个正则方程 \(b,b'\)。它们在原厚化上的理想相同，故可用一个提升的单位重新归一化，使两方程在模 \(V\) 后相同。两方程之差落在 \(V\otimes_k\OO_{X_0}\)，取其模约化方程 \(b_0\) 的类。这个类定义了从 \(\mathcal I_{D_0}/\mathcal I_{D_0}^2\) 到 \(V\otimes_k\OO_{D_0}\) 的线性映射：换方程只按方程的单位过渡系数改变，换归一化单位则多出一个 \(b_0\) 的倍数。这里用 \(X\) 的平坦性识别增量理想，并用相对有效除子换底后仍是有效 Cartier 除子（\href{https://stacks.math.columbia.edu/tag/056P}{Stacks，056P}）。如果这个类为零，方程相差的增量是约化方程的倍数，可由一个形如 \(1+c\)、\(c\) 落在增量理想中的单位吸收，故两除子相同。局部映射粘合，得到 \(H^0(D_0,\mathcal N_{D_0/X_0})\otimes_kV\) 中的差。

对于态射 \(X\to P\)，可在同像点附近取靶的仿射邻域，比较两环映射；由同样的导子规则，差粘合为 \(\Hom_{\OO_{X_0}}(f_0^*\Omega^1_{P_0/k},\OO_{X_0})\otimes_kV\) 中的元素。此处不要求该全局 Hom 是有限维的；只要 \(V\) 有限维，取一组基便能把相关张量写成有限直和。在对偶数环上得到的是切空间映射；任意此类 \(V\) 的差映射按同样计算为其与 \(V\) 的张量。检查切空间的单射须用纤维函子在对偶数环上的单射，不能改用域值点。b) 的函子态射此处按4.1解释为 \(u\) 的逆转置与基点补开上 \(f\) 的复合。

\paragraph{a) 的任意底：只用像点的纤维。}
在仿射底上，取有限生成的整数子代数作诺特模型。有限表示的 \(X\)、参数丛、两截面、通用除子及两截面所给除子的相等都可在某一有限阶段实现：可先降下闭浸入、线丛及其截面，再用有限表示态射的极限性质使闭子概形的相等成立。由平坦性下降，可使 \(X\) 在模型底上及通用除子在模型参数丛上平坦（\href{https://stacks.math.columbia.edu/tag/01ZL}{Stacks，01ZL}；\href{https://stacks.math.columbia.edu/tag/081C}{Stacks，081C}）。通用除子的线丛方程也可在有限次放大后保持正则：局部乘法映射的余核已对参数底平坦，且源、靶均平坦。于是其核随底扩张而换底；模型上的核有限生成，在极限为零，有限次放大即可消去。故有模型上的相对有效 Cartier 除子，并可用它定义比较函子。

设 \(t\) 是模型底上一个来自实际底点 \(s\) 的像点。模型比较态射在 \(k(t)\) 上的纤维，经扩域到 \(k(s)\) 后，是已完成域底论证的比较态射。由忠实平坦性，它在 \(k(t)\) 上已是单态射：任取同像的两检验截面，扩域后相等，而射影丛的截面相等可沿这种覆盖检验。这不要求模型的所有纤维都满足原稿关于关联循环的几何条件。仅在 \(t\) 处的诺特局部环上作逐阶归纳，只会用这个闭点的纤维；克鲁尔交定理给出两个模型截面在该局部环上相等。拉回至 \(\OO_{S,s}\)，对每个 \(s\) 应用，得所需的相等。

\paragraph{b) 的任意底：可直接检测标架。}
在秩为 \(r>0\) 的部分，令 \(U=X-Z\)。在 \(U\) 的 \((r+1)\) 重纤维积中，像在 \(P\) 里构成射影标架的条件是开条件；局部由各 \(r\) 列行列式为单位给出。这个开子概形在每条几何纤维均非空，由域底标架引理可知。它平坦、局部有限表示、满射于底，所以给出忠实平坦覆盖。在该覆盖上有通用标架截面；若两自同构与 \(f\) 的复合相等，它们在该标架上的作用相同，因此自同构相等。最后沿覆盖检测相等即可。这是 b) 的单独证明，不能直接用于 a) 的参数截面单态射。
\end{document}
